\ncdlesson{2}{Skal vi lige tage en kop kaffe?}{我们要不喝杯咖啡？}
{DU3 Modul 4（B1+ $\rightarrow$ B2-）｜对标 NCE 3 \& Modultest 4 \textit{Læsning Opgave 2 (Chat) \& Mundtlig Opgave 2 (Paropgave)}}
{\textbf{灵魂情态小品词 II}：\textit{lige}（稍微/顺便/降压）· \textit{vel}（反意求同“吧？”）· \textit{bare}（只是/尽管）· \textit{egentlig}（究竟/其实）· \textit{ellers}（要不然/除此之外）的语气软化与礼貌协商功能}
{在丹麦职场与日常对话中，直译中文的祈使句或英语的 \textit{Can you help me?} 常显得过于生硬，甚至带着命令感。学会用 \textbf{\textit{lige, vel, bare, egentlig, ellers}} 给请求“卸下压迫感”，是融入丹麦平权文化与通过 DU3 Modul 4 口语双人互动的关键。本课带你听懂咖啡机旁的真实协商艺术！}

% ---------------------------------------------------------------------
\begin{foerduleaeser}
带着下面两个问题阅读 Tekst A 和 Tekst B：
\begin{enumerate}
  \item \textit{Hvad vil Kasper have Emma til at hjælpe med i Tekst A, og hvordan giver han hende plads til at sige nej?}（在 Tekst A 中，卡斯珀想请艾玛帮什么忙？他是如何给对方留出拒绝余地的？）
  \item \textit{Hvorfor tøver Sofie med at tage med til fredagens borgermøde i Tekst B, og hvilket alternativ foreslår Thomas med ”ellers”?}（在 Tekst B 中，索菲为什么对参加周五的市民会议有些犹豫？托马斯用 \textit{ellers} 提出了什么替代方案？）
\end{enumerate}
\end{foerduleaeser}

% ---------------------------------------------------------------------
\ncdsubsec{二、Tekster（口语微调 / 报刊与思辨双语篇）}

\subsection*{Tekst A：\textit{Lige et øjeblik over kaffen}（职场口语对话 · 约 205 词）}

\begin{tipbox}{语篇说明}
本篇还原丹麦办公室咖啡机旁的真实对话。阅读时请特别关注\textbf{\color{NCDNavy}深蓝色加粗的 \textit{lige / bare / egentlig}}（用于放低姿态、发出邀请或温和切入）与\textbf{\color{NCDRed}红色加粗的 \textit{vel / ellers}}（用于寻求认同与提供退路）。
\end{tipbox}

\begin{tcolorbox}[enhanced, breakable, colback=NCDBgBlue, colframe=NCDNavy, coltitle=white, fonttitle=\bfseries, title={办公室咖啡间协商｜Kasper \& Emma}, arc=2pt, boxrule=0.8pt, left=10pt, right=10pt, top=7pt, bottom=7pt]
\small
\textbf{Kasper}: Hej Emma! Har du \da{egentlig} travlt lige nu, eller skal vi \da{lige} tage en kop kaffe?\\[3pt]
\textbf{Emma}: Hej Kasper! Jeg har \da{bare} et par mails tilbage, men en kaffepause lyder dejlig. Hvad tænker du på?\\[3pt]
\textbf{Kasper}: Jeg skal fremlægge vores nye spareplan for bestyrelsen i morgen. Kunne du \da{lige} kigge mine fem slides igennem bagefter? Det tager \dabold{vel} kun ti minutter?\\[3pt]
\textbf{Emma}: Jo, selvfølgelig vil jeg det! Men hvordan går det \dabold{ellers} med tallene? Du virkede lidt presset i går.\\[3pt]
\textbf{Kasper}: Ja, det var lidt stramt, men nu ser det fint ud. Du må \da{bare} sige til, hvis du hellere vil gøre det efter frokost. \dabold{Ellers} kan jeg også sende udkastet til dig på mail, så du \da{bare} læser det, når det passer dig.\\[3pt]
\textbf{Emma}: Nej, lad os \da{bare} tage det klokken tolv ved mit skrivebord. Vi skal \da{lige} sikre, at konklusionen står skarpt, ikke?\\[3pt]
\textbf{Kasper}: Jo, præcis! Det er \dabold{vel} heller ikke nødvendigt at tage alle bilagene med i selve præsentationen?\\[3pt]
\textbf{Emma}: Nej, det er det bestemt ikke. Hent du kaffen, så rydder jeg \da{lige} mit bord imens!
\end{tcolorbox}

\begin{transbox}
\small
\textbf{\color{NCDTeal}[Tekst A 中文精译]}\\
\textbf{卡斯珀}：嗨，艾玛！你现在其实忙吗，还是我们要不要顺便喝杯咖啡？\\
\textbf{艾玛}：嗨，卡斯珀！我只剩几封邮件要处理，不过去喝杯咖啡休息一下听起来很棒。你在想什么事呢？\\
\textbf{卡斯珀}：我明天得向理事会汇报我们的新节约方案。待会儿你能稍微帮我看一下我那五页幻灯片吗？大概也就花十分钟吧？\\
\textbf{艾玛}：好呀，我当然愿意！不过除此之外，数据那块推进得怎么样？你昨天看起来压力有点大。\\
\textbf{卡斯珀}：是啊，当时有点赶，不过现在看上去不错了。如果你更愿意午饭后再看，尽管告诉我。要不然我也可以把草稿邮件发给你，等你方便的时候再看就行。\\
\textbf{艾玛}：不用，我们十二点就在我办公桌那儿直接看吧。我们稍微确认一下结论写得够犀利，对吧？\\
\textbf{卡斯珀}：对，正是！反正也没必要把所有附件都放进正式演示里吧？\\
\textbf{艾玛}：是的，完全没必要。你去倒咖啡吧，我顺便把桌子收拾一下！
\end{transbox}

\subsection*{Tekst B：\textit{Chat om fredagens borgermøde}（对标 DU3 Modul 4 社交短信协商 · 约 140 词）}

\begin{emailbox}{微信 / 短信对话｜Sofie \& Thomas}
\textbf{Thomas}: Hej Sofie! Skal vi ikke tage sammen hen til borgermødet om nye cykelstier på fredag?\\[3pt]
\textbf{Sofie}: Hej Thomas! Jeg vil \da{egentlig} gerne med, men mødet slutter \dabold{vel} ret sent, gør det ikke? Jeg skal nemlig tidligt op lørdag morgen.\\[3pt]
\textbf{Thomas}: Det varer vist til omkring klokken ni, men du kan \da{bare} gå i pausen efter første time, hvis du bliver træt. \dabold{Ellers} kan vi \da{lige} mødes til en hurtig kop kaffe klokken fem først, så vi når at drøfte vores spørgsmål inden mødet.\\[3pt]
\textbf{Sofie}: Åh, det er en super idé! Jeg tjekker \da{lige} min kalender, når jeg kommer hjem, og så skriver jeg til dig i aften. Vi kan \da{vel} mødes ved biblioteket?\\[3pt]
\textbf{Thomas}: Ja da! Sig \da{bare} til, hvis klokken fem bliver for tidligt for dig.
\end{emailbox}

\textbf{Tekst B 阅读核对（Rigtigt / Forkert / Kort svar，答案见第六部分）：}
\begin{enumerate}
  \item \textit{Sofie vil slet ikke med til borgermødet om cykelstier, fordi emnet ikke interesserer hende.} \underline{\hspace{2.2cm}}
  \item \textit{Hvad foreslår Thomas, at Sofie kan gøre, hvis mødet bliver for langt for hende?} \\\underline{\hspace{5.2cm}}
  \item \textit{Hvornår lover Sofie at vende tilbage med endeligt svar til Thomas?} \underline{\hspace{4.5cm}}
\end{enumerate}

\begin{tipbox}{本课发音与语调要点}
\begin{itemize}
  \item \textbf{五大软化小品词的口语极简吞音}：\textit{lige} 读 \ipa{ˈliːə}（\textbf{g 完全不发音}，快读时常缩成短促的 \ipa{li}）；\textit{egentlig} 读 \ipa{ˈeːˀənli}（\textbf{g 与 t 双静音}！）；\textit{bare} 读 \ipa{ˈbɑːɐ}（常弱化为 \ipa{bɑ}）；\textit{vel} 在句中轻读 \ipa{vɛl}，但在句末疑问标签 \textit{..., vel?} 中音调上扬；\textit{ellers} \ipa{ˈɛlˀɐs}（带清晰喉塞音）。
  \item \textbf{对话连读与弱读节奏}：\textit{Skal vi lige} 常连读为 \ipa{sga vi ˈliːə}（\textit{skal} 的 l 常脱落）；\textit{Det tager vel} \ipa{de ˈtˢɑːˀ vɛl}（\textit{tager} 的 g 不发音！）。
  \item \textbf{核心词重音与软 d}：\textit{et borgermøde} \ipa{ˈbɒʊ̯ˀɐˌmøːðə}（\textit{borger} 中的 \textit{g} 滑音化为 \ipa{ʊ̯}，\textit{møde} 为软 d）；\textit{behagelig} \ipa{beˈhæːˀəli}；\textit{et udkast} \ipa{ˈuðˌkʰasd}。
\end{itemize}
\end{tipbox}

\begin{transbox}
\small
\textbf{\color{NCDTeal}[Tekst B 中文对照]}\\
\textbf{托马斯}：嗨，索菲！这周五我们要不要一起去参加关于新建自行车道的市民会议？\\
\textbf{索菲}：嗨，托马斯！我其实很想去，不过会议大概会结束得挺晚的，对吧？我周六一大早还要早起呢。\\
\textbf{托马斯}：听说会持续到九点左右，但如果你觉得累，完全可以在第一小时结束后的休息时间先走呀。要不然我们也可以五点先稍微碰头喝杯快咖啡，这样赶在开会前就能讨论好我们的问题。\\
\textbf{索菲}：啊，这主意太棒了！我回家后稍微查一下日程表，今晚给你回信。我们大概约在图书馆门口碰面吧？\\
\textbf{托马斯}：那当然！如果五点对你来说太早，尽管告诉我一声就行。
\end{transbox}

% ---------------------------------------------------------------------
\ncdsubsec{三、Gloser（核心生词表）}

\begin{tipbox}{词表重点}
本课五个核心小品词 \textit{lige, vel, bare, egentlig, ellers} 在句中作语气词时都没有词形变化；但务必注意 \textit{at gennemgå}（浏览、审核）和 \textit{at fremlægge}（汇报、呈交）是不规则复合动词，过去时和完成时分别随 \textit{gå} 和 \textit{lægge} 变化！
\end{tipbox}

\begin{glosertable}
小品词 & lige & \textit{(不变，情态小品词)} & 稍微，顺便，一下子（软化请求） \\
小品词 & vel & \textit{(不变，情态小品词)} & 大概……吧，难道不是吗（求同） \\
小品词 & bare & \textit{(不变，情态小品词)} & 只是，尽管，放心……好了 \\
小品词 & egentlig & \textit{(不变，情态小品词)} & 其实，究竟，说起来 \\
副词 & ellers & \textit{(不变，副词/小品词)} & 要不然，否则；除此之外 \\
\midrule
名词 & et øjeblik & \textit{(øjeblikket, øjeblikke, -ne)} & 片刻，一会儿 \\
名词 & en kaffepause & \textit{(kaffepausen, -pauser, -erne)} & 喝咖啡小憩时间 \\
名词 & en spareplan & \textit{(spareplanen, -planer, -erne)} & 节支方案，紧缩计划 \\
名词 & en bestyrelse & \textit{(bestyrelsen, bestyrelser, -ne)} & 理事会，董事会 \\
名词 & et borgermøde & \textit{(borgermødet, -møder, -erne)} & 市民议事会，居民大会 \\
名词 & et bilag & \textit{(bilaget, bilag, bilagene)} & 附件，附录（单复同形） \\
名词 & en konklusion & \textit{(konklusionen, -er, -erne)} & 结论 \\
\midrule
动词 [har] & at fremlægge & \textit{(fremlægger, fremlagde, har fremlagt)} & 汇报，展示，提交（方案） \\
动词 [har] & at gennemgå & \textit{(gennemgår, gennemgik, har gennemgået)} & 过一遍，仔细审核 \\
动词 [har] & at drøfte & \textit{(drøfter, drøftede, har drøftet)} & 探讨，商议（偏书面） \\
动词 [har] & at sikre & \textit{(sikrer, sikrede, har sikret)} & 确保，保障 \\
\midrule
形容词 & krævende & \textit{(不变，现在分词形容词)} & 要求苛刻的，费心力的 \\
\end{glosertable}

% ---------------------------------------------------------------------
\ncdsubsec{四、Grammatik \& Stilistik（核心语法、语气与修辞精讲）}

\begin{grammatikbox}{语法点 1：五大语气软化小品词（\textit{lige, vel, bare, egentlig, ellers}）语义与语用矩阵}
如果说 Lektion 1 的 \textit{jo, da, nok, vist} 负责交代“信息真伪与双方共识”，那么本课的五个小品词就是丹麦社交润滑剂——专门用来\textbf{降低指令侵略性、缓和冲突与制造亲切感}：

\vspace{0.3em}
\renewcommand{\arraystretch}{1.25}
{\small
\begin{tabularx}{\linewidth}{@{} >{\bfseries\raggedright\arraybackslash}p{1.65cm} >{\raggedright\arraybackslash}p{3.4cm} >{\raggedright\arraybackslash}p{4.3cm} >{\raggedright\arraybackslash}X @{}}
\toprule
\textbf{小品词} & \textbf{核心降压功能} & \textbf{语用潜台词} & \textbf{地道高频句式} \\
\midrule
\da{lige} & 弱化请求分量\newline （“稍微 / 顺便”） & “就花你一小会儿，不是大事” & \textit{Kan du \da{lige} åbne vinduet?}\newline \textit{Jeg tjekker \da{lige} min kalender.} \\
\midrule
\dabold{vel} & 礼貌寻求确认\newline （“大概……吧？”） & “我猜是这样，对不对？”（暗含期待对方点头） & \textit{Det koster \dabold{vel} ikke meget?}\newline \textit{Du kommer med, \dabold{vel}?}（句末标签） \\
\midrule
\da{bare} & 授权、安慰或轻描淡写（“尽管 / 只是”） & “没关系，随便这么做就好 / 别多想” & \textit{Tag \da{bare} den sidste bolle!}\newline \textit{Du må \da{bare} ringe, hvis der er noget.} \\
\midrule
\da{egentlig} & 缓冲突兀提问或委婉转折（“究竟 / 其实”） & “顺便好奇问一句 / 说实在的……” & \textit{Hvad synes du \da{egentlig} om planen?}\newline \textit{Jeg vil \da{egentlig} gerne med, men...} \\
\midrule
\dabold{ellers} & 留出体面退路或关切寒暄（“要不然 / 另外”） & “如果不方便也没事，我们还有 B 方案” & \textit{\dabold{Ellers} kan vi mødes i morgen.}\newline \textit{Hvordan går det \dabold{ellers} på jobbet?} \\
\bottomrule
\end{tabularx}
}
\vspace{0.3em}

\textbf{避坑对比}：\ding{172} 询问对方意愿时，说 \textit{Hvad mener du?} 听起来像审讯；加上 \textit{egentlig}（\textit{Hvad mener du \da{egentlig}?}）立刻变成真诚好奇的探讨。\ding{173} \textit{ellers} 放句首表示“否则/要不然”时，\textbf{触发 V2 倒装}：\textit{\dabold{Ellers} \textbf{kan} vi tage det i morgen.}（绝不能写成 \ding{55}\ \textit{Ellers vi kan...}）。
\end{grammatikbox}

\begin{grammatikbox}{语法点 2：高情商请求的“三步软化法”与句末 \textit{, vel?} 的回答规则}
\begin{enumerate}
  \item \textbf{从生硬祈使句升级到丹麦礼貌体}：
    \begin{itemize}
      \item 生硬直译：\ding{55}\ \textit{Hjælp mig med opgaven!}（带有“命令”感，只适合极其紧急情况）
      \item 柔和邀请：\ding{51}\ \textit{Kunne du \da{lige} hjælpe mig med opgaven, hvis du har tid?}（过去式情态动词 \textit{Kunne} + \textit{lige} 双重软化）
      \item 体贴留退路：\ding{51}\ \textit{Du kan \da{bare} kigge på den i morgen --- \dabold{ellers} finder vi en anden løsning.}
    \end{itemize}
  \item \textbf{句末确认标签 \textit{, vel?} 与 \textit{, ikke?} 的细微区别及回答法}：
    \ding{172} \textit{Det er dyrt, \textbf{ikke}?}（说话人几乎确定，只等对方附和）；\textit{Det er ikke for dyrt, \dabold{vel}?}（说话人希望得到“不贵”的安抚，带有小心翼翼试探感）。\newline
    \ding{173} \textbf{回答否定式 \textit{vel?} 疑问句注意 \textit{Jo} vs \textit{Nej}}：若对方问 \textit{Det tager ikke lang tid, vel?}（不会花很久吧？），如果你同意“确实不花很久”，请回答 \da{Nej, slet ikke}；如果你想反驳“不，其实要花挺久”，则用重读的 \dabold{Jo, det tager faktisk en time!}
\end{enumerate}
\end{grammatikbox}

% ---------------------------------------------------------------------
\ncdsubsec{五、Mønsterøvelser（句型与修辞操练）}

\textbf{练习 A：高情商软化改写——按括号提示用指定小品词软化下列生硬句子（注意语序与 V2 倒装）}
\begin{enumerate}
  \item \textit{Send udkastet til mig i dag.}（改用 \textit{Kan du lige...} 软化为礼貌请求）$\rightarrow$
  \item \textit{Hvad arbejder du med nu?}（在动词和主语后插入 \textit{egentlig}，并改得更柔和）$\rightarrow$
  \item \textit{Hvis du ikke har tid, finder vi en ny dag.}（以 \textit{Ellers} 开头改写后半句，注意 V2 倒装）$\rightarrow$
  \item \textit{Du må bruge min computer.}（在祈使句或情态句中插入 \textit{bare} 表热情许可）$\rightarrow$
  \item \textit{Du har ikke glemt mødet i morgen?}（用句中 \textit{vel} 改写为委婉确认问句）$\rightarrow$
  \item \textit{Hvordan går det med dit nye arbejde?}（插入 \textit{ellers} 表示“除此之外最近如何”）$\rightarrow$
\end{enumerate}

\textbf{练习 B：语境选词填空（从 \textit{lige, vel, bare, egentlig, ellers} 中选填，注意第 6 题双空！）}
\begin{enumerate}
  \item Undskyld, kan jeg \underline{\hspace{2.2cm}} låne din kuglepen et øjeblik?
  \item Du skal ikke bekymre dig om regningen --- jeg skal \underline{\hspace{2.2cm}} betale kaffen i dag!（表“尽管让我来就行”）
  \item Det bliver \underline{\hspace{2.2cm}} ikke regnvejr i eftermiddag, gør det?
  \item Jeg troede, at han boede i Odense, men han kommer \underline{\hspace{2.2cm}} fra Esbjerg.
  \item Lad os mødes klokken tre, og \underline{\hspace{2.2cm}} kan vi altid ringes ved i aften.
  \item Vil du \underline{\hspace{2.2cm}} holde døren et sekund? Du behøver \underline{\hspace{2.2cm}} ikke at vente på mig bagefter.
\end{enumerate}

\textbf{练习 C：中译丹实战（自然融入五大语气软化词）}
\begin{enumerate}
  \item 我们要不要顺便去咖啡馆喝杯咖啡？（用 \textit{lige}）
  \item 说起来，你周末究竟打算做什么？（用 \textit{egentlig}）
  \item 如果你有问题，尽管给我发消息就行。（用 \textit{bare}）
  \item 这次考试大概不会太难吧？（用 \textit{vel}）
  \item 我们可以今天见面，要不然就周四再聊。（用 \textit{ellers}）
  \item 顺便问一下，除此之外你的项目进展得怎么样？（用 \textit{ellers}）
\end{enumerate}

% ---------------------------------------------------------------------
\ncdsubsec{六、DU3 Modul 4 实战微任务与参考答案}

\begin{emailbox}{DU3 Modul 4 写作任务（对标 Skrivning Opgave 1 · 私人协商邮件 $\ge$ 125 ord）}
\textbf{Situation}: Du skal holde et vigtigt oplæg på dit studie eller arbejde næste tirsdag, men du er usikker på, om din tekst er klar og naturlig. Din danske kollega eller studieven, Sofie, er altid meget hjælpsom, men hun har også meget travlt for tiden.\\
\textbf{Opgave}: Skriv en høflig og uformel mail til Sofie på \textbf{mindst 125 ord} (gerne 140--160 ord). Spørg forsigtigt, om hun har tid til at læse dit udkast igennem eller mødes over en kop kaffe, giv hende en nem mulighed for at sige nej, og brug alle fem ord (\textit{lige, vel, bare, egentlig, ellers}) naturligt i mailen!
\end{emailbox}

\begin{answerbox}
\textbf{Før du læser 参考回答}
\begin{enumerate}
  \item \textit{Kasper vil gerne have Emma til at kigge hans fem slides igennem inden bestyrelsesmødet. Han giver hende plads til at sige nej ved at sige ”Du må bare sige til, hvis du hellere vil gøre det efter frokost” og foreslå at sende udkastet på mail i stedet.}
  \item \textit{Sofie tøver, fordi mødet slutter ret sent, og hun skal tidligt op lørdag morgen. Med ”ellers” foreslår Thomas, at de kan mødes til en hurtig kop kaffe klokken fem først for at drøfte deres spørgsmål inden mødet.}
\end{enumerate}

\textbf{Tekst B 阅读核对答案}
\begin{enumerate}
  \item \da{Forkert}（索菲表示她其实很想去：\textit{Jeg vil egentlig gerne med}，只是担心周五结束太晚影响周六早起）
  \item \da{Han foreslår, at hun bare kan gå i pausen efter første time, hvis hun bliver træt.}
  \item \da{I aften, når hun har tjekket sin kalender derhjemme.}
\end{enumerate}

\textbf{练习 A 答案}
\begin{enumerate}
  \item \textit{Kan du \da{lige} sende udkastet til mig i dag?}（用 \textit{Kan du lige...} 彻底消除命令感）
  \item \textit{Hvad arbejder du \da{egentlig} med nu?}（在主语 \textit{du} 之后插入 \textit{egentlig}）
  \item \textit{\dabold{Ellers} finder vi bare en ny dag.}（\textit{Ellers} 放句首后动词 \textit{finder} 紧跟占第 2 位）
  \item \textit{Du må \da{bare} bruge min computer.}（或 \textit{Brug bare min computer!}，\textit{bare} 传达热情大方）
  \item \textit{Du har \dabold{vel} ikke glemt mødet i morgen?}（主句中 \textit{vel} 放在 \textit{ikke} 前面，句末语调上扬）
  \item \textit{Hvordan går det \dabold{ellers} med dit nye arbejde?}（位于主语 \textit{det} 之后表顺便关切）
\end{enumerate}

\textbf{练习 B 答案}
\begin{enumerate}
  \item \da{lige}（“稍微借用一下”，降低打扰感）
  \item \da{bare}（“放宽心，由我来付吧”）
  \item \dabold{vel}（反意疑问句寻求安慰确认：“下午大概不会下雨吧？”）
  \item \da{egentlig}（点出事实背后的真实情况：“其实来自埃斯比约”）
  \item \dabold{ellers}（提供替代方案：“要不然我们也可以晚上通电话”）
  \item \da{lige}；\da{bare}（第一空用 \textit{lige} 请求顺手扶一下门；第二空用 \textit{bare} 宽慰对方“你完全不用等我”）
\end{enumerate}

\textbf{练习 C 答案}
\begin{enumerate}
  \item \textit{Skal vi lige tage en kop kaffe på caféen?}
  \item \textit{Hvad skal du egentlig lave i weekenden?}
  \item \textit{Hvis du har spørgsmål, kan du bare sende mig en besked.}
  \item \textit{Denne eksamen bliver vel ikke alt for svær?}
  \item \textit{Vi kan mødes i dag; ellers snakker vi bare på torsdag.}
  \item \textit{Hvordan går det ellers med dit projekt?}
\end{enumerate}

\textbf{DU3 Modul 4 高分范文（146 词，不含署名）}\\
\textit{Hej Sofie}

\textit{Hvordan går det \dabold{ellers} med dig og dit nye projekt? Jeg håber, at du ikke har alt for meget om ørerne i denne uge.}

\textit{Jeg skriver \da{egentlig}, fordi jeg skal holde et vigtigt oplæg næste tirsdag, og jeg er stadig lidt usikker på sproget i min konklusion. Jeg ved godt, at du har travlt, men kunne du \da{lige} kigge mit korte udkast igennem? Det tager \dabold{vel} kun ti minutter, og vi kunne måske tage en hurtig kop kaffe i kantinen efter frokost?}

\textit{Hvis din kalender er helt fyldt i morgen, må du \da{bare} sige nej --- det forstår jeg selvfølgelig fuldt ud! \dabold{Ellers} kan jeg også sende teksten på mail, så du \da{bare} læser den, hvis du får et roligt øjeblik før weekenden.}

\textit{Sig \da{bare} til, hvad der passer dig bedst!}

\textit{De bedste hilsner\\
\mbox{[Dit navn]}}

\small\textbf{名师点评}：146 词（不含署名）流畅自然地达到了 Modul 4 私人协商邮件的 B1+/B2- 水准。开篇用 \textit{ellers} 寒暄切入，中段用 \textit{egentlig} 温和引出请求、\textit{lige} 与 \textit{vel} 降低对方负担，最后用 \textit{bare} 和句首倒装的 \textit{Ellers kan jeg...} 留出充分退路，既展现了对情态小品词的纯熟驾驭，又极其贴合丹麦职场平等体贴的沟通礼仪。
\end{answerbox}